\section*{Техническое задание}

\begin{table}[H]
    \centering
    \caption{Исходные данные для варианта 66}
    \label{table:params}
    \begin{tblr}{
        colspec = {X[c]X[c]X[c]X[c]X[c]},
        hlines,
        vlines,
        rows = {abovesep=8pt,belowsep=8pt},
    }
        Схема ОС & $b$ & $c$ & $J_{д}$ \\
        $k_{\text{пр}},\,k_{\text{им}}$ & 4 & $10^3$ & $2 \cdot 10^{-3}$ \\
    \end{tblr}
\end{table}

Выбрать параметры аналогового регулятора модуля поворота промышленного робота.

Углы поворота ротора двигателя $\varphi_{р}$ и выходного вала редуктора $\varphi_{м}$ измеряются и пропорциональные им сигналы $q_{р}$ и $q_{м}$ подаются на блок сравнения с задающим воздействием $q^* = k\varphi^*$, где $k$ — коэффициент преобразования измерителя угла поворота, $i$ — передаточное отношение редуктора, $\varphi^*$ — требуемый закон перемещения модуля поворота.

Сигналы рассогласования $\psi_{р} = q_{р} - q^*$ и $\psi_{м} = q_{м} - q^*$ поступают на аналоговый регулятор. В соответствии с вариантом используется \textit{\textbf{ПИ-регулятор}}, содержащий интегральное звено по сигналу $\psi_{р}$ с коэффициентом $k_{и}$ и пропорциональное звено по сигналу $\psi_{м}$ с коэффициентом $k_{п}$.

Сформированный отрицательной обратной связью сигнал управления $u$ подается на вход двигателя (рисунок~1).

% ВСТАВИТЬ ЗДЕСЬ СКРИНШОТ СТРУКТУРНОЙ СХЕМЫ
% \begin{figure}[H]
%     \centering
%     \includegraphics[width=0.8\textwidth]{img/your_scheme.png}
%     \caption{Структурная схема системы управления}
%     \label{fig:scheme}
% \end{figure}

Момент инерции приводимого в движение выходного звена $J_{м} = 1$ кг$\cdot$м$^2$. Жесткость и коэффициент демпфирования редуктора, приведенные к выходному валу: $c = 10^3$ Нм, $b = 4$ Нм$\cdot$с. Момент инерции ротора двигателя $J_{д} = 2 \cdot 10^{-3}$ кг$\cdot$м$^2$.

При расчетах необходимо использовать линейную динамическую характеристику двигателя
\[
\tau \dot M_{\text{д}} + M_{\text{д}} = -s \dot \varphi_{р} + r u,
\]
где $s$ — крутизна механической характеристики, $\tau$ — собственная постоянная времени двигателя, $r$ — коэффициент пропорциональности, $M_д$ — движущий момент.

Записать уравнения движения механической части как системы с двумя степенями свободы, приведенную динамическую характеристику двигателя и уравнение системы управления. Разрешить полученную систему четырех уравнений в переменных Лапласа относительно четырех неизвестных. Сформировать знаменатель передаточных функций.

Определить область устойчивости замкнутой системы управления с помощью критериев Стодолы и Рауса–Гурвица. На диаграмме с $D$-разбиением при $\varphi^* \equiv 180^\circ$ построить область крутящего момента $\{M\}$, в которой $|M_{м}| < 150$ Нм, и область управляющего сигнала $\{U\}$, в которой $|u| < 220$ В.

В области одновременного выполнения этих условий построить линии равных длительностей переходных процессов $t_{п} = t_{п}(k_{п}, k_{и})$ при нулевых начальных условиях. Под длительностью переходного процесса понимать время, после которого ошибка $|\psi_{м}|$ не превышает 10\% от $\varphi^*$.

Построить оптимальную кривую по функционалу качества
\[
J = \int_0^{t_1} \left(\psi^2 + \rho u^2\right)\,dt.
\]

Найти значения параметров $k_{п}$ и $k_{и}$, соответствующие наименьшей длительности переходных процессов. Проиллюстрировать результаты расчетов при нескольких значениях $(k_{п}; k_{и})$, построив графики перемещения $\varphi_{м}$, сигнала рассогласования $\psi_{р}$, упругой деформации редуктора $\varphi_{р}/i - \varphi_{м}$, крутящего момента $M_{м}$, приведенного момента двигателя и сигнала управления $u$ как функций времени.
